\chapter{Conclusion}\label{ch:conclusion}

\section{Answers to the research questions}
\begin{description}[leftmargin=1.2cm,style=nextline]
  \item[RQ1] The \NWaves{} waves are labelled by a descriptive rule, and PELT placed \NChangepoints{} change points in the weekly series (\NChangepointsLo{} to \NChangepointsHi{} under other settings). Both describe the curve and neither is an estimate. In weekly first differences, a one-point rise in stringency went with a rise of \PolSlopeA{} in the mobility retreat index in 2020, \PolSlopeB{} in 2021 and \PolSlopeC{} in 2022 (2022 against 2020, \PolIntCp). The 2020 slope is about \PolRatio{} times the 2022 slope, and the contrast holds without March 2020 ($p$~\PolIntCpNoMarch).
  \item[RQ2] Mean diet diversity rose by \PanFgDiff{} food groups (95\% CI \PanFgCIlo{} to \PanFgCIhi) and MDD-W attainment from \PanMddBefore\% to \PanMddAfter\%, while \PanPPP\% of women reported a purchasing-power constraint. I found no evidence that the gain differed by wealth. It was \PanGainQOne{} food groups in the poorest quintile and \PanGainQFive{} in the richest (joint test, \PanChgP). The wealth gap in MDD-W was present in 2019 (\PanQOneMddB\% against \PanQFiveMddB\%) and in 2020 (\PanQOneMddA\% against \PanQFiveMddA\%). A change from in-person to telephone interviewing could also produce a uniform shift, and these data cannot separate the two explanations.
  \item[RQ3] The published prevalences were reproduced exactly, at \PrevDep\% (95\% CI \PrevDepLo--\PrevDepHi) for probable depression and \PrevAnx\% for probable anxiety. After adjustment, depression was more common among respondents who had lost a family member to COVID-19 (aOR \OrFamDeathB, aPR \PrFamDeathB), reported a prior mental-health problem (aPR \PrHistB) or were students (aPR \PrStudentB), and less common among those with two vaccine doses (aPR \PrTwoDoseB). Chattogram had the lowest prevalence and \SuspectChatto{} of the \SuspectN{} flagged records, none of them a depression case. Without those records its adjusted odds ratio for depression moved from \SensChattoAll{} to \SensChattoExcl.
  \item[RQ4] ICL selected \GmmK{} profiles (entropy \GmmEntropy), but the mean silhouette is \GmmSil{} and BIC is highest at the edge of the grid. Profile membership explains \ProfPcOneRsq{} of the variance in the first principal component of the indicators, so the profiles are best described as ordered bands of one severity score. On the full sample, a small class dominated by flagged records appeared in \SpikeFits{} of the \SpikeFitsTotal{} estimable diagonal fits.
  \item[RQ5] Over \RepN{} repeated splits, logistic regression reached a mean test AUC of \RepAucALogit{} from sociodemographic and COVID-19 variables, \RepAucBLogit{} after adding optimism, pessimism and mindfulness, and \RepAucCLogit{} after adding anxiety and stress scores. No learner's mean AUC differed from it by more than \RepMaxAbsDiff, against a split-to-split standard deviation of \RepSdMin--\RepSdMax. Mindfulness accounts for most of the tier-B gain (\AblAMind{} when added to tier A alone, against \AblAOpt{} for optimism and \AblAPess{} for pessimism). Two PHQ-9 items alone reached \AblPhqTwo, a ceiling that overlaps the outcome. The psychological scales are concurrent self-reports, so these figures measure shared information and do not support a claim about screening.
\end{description}

The published prevalences reproduce, and the main associations keep their direction when the flagged records are excluded. The evidence comes from an online convenience sample, a two-district panel and national aggregates, and should not be generalised to the population of Bangladesh.

\section{Future work}
The individual-level analysis should be repeated on a probability sample with repeated measures, which would give population prevalence and within-person change, and the prediction results should be tested on a survey collected by other investigators. For the diet panel, a mode experiment that randomises women to in-person or telephone recall would give the direction and size of the mode effect. The screen for templated response blocks could be run on other open online surveys from the period to see how common such blocks are. The PSS-4 also needs a psychometric study in Bangladeshi samples, given its low reliability here ($\alpha = \AlphaPSS$).
